% !TEX program = xelatex
% Lernplatt - by Can Siebert
% Kurs 71 (Dig 42) - Tag 7 - Lehrskript
% Process Mining: Event Logs -> Discovery -> Conformance als Steuerungsinstrument
%
% Self-contained xelatex source. Kein biblatex, keine externen Bilder.

\documentclass[11pt,a4paper,oneside]{article}

% --- Encoding & Sprache --------------------------------------------------
\usepackage{polyglossia}
\setdefaultlanguage{german}

% --- Geometrie -----------------------------------------------------------
\usepackage[a4paper,top=2cm,bottom=2cm,outer=2.5cm,inner=2.5cm,bindingoffset=1.5cm]{geometry}

% --- Fonts ---------------------------------------------------------------
\usepackage{fontspec}
\defaultfontfeatures{Ligatures=TeX}

\IfFontExistsTF{Charter}{%
  \setmainfont{Charter}[%
    Numbers=OldStyle,%
    UprightFeatures   = {SmallCapsFeatures={Letters=SmallCaps}}%
  ]%
}{%
  \IfFontExistsTF{Bitstream Charter}{%
    \setmainfont{Bitstream Charter}%
  }{%
    \setmainfont{TeX Gyre Schola}%
  }%
}

\IfFontExistsTF{Source Sans 3}{%
  \setsansfont{Source Sans 3}%
}{%
  \IfFontExistsTF{Source Sans Pro}{%
    \setsansfont{Source Sans Pro}%
  }{%
    \setsansfont{TeX Gyre Heros}%
  }%
}

\newcommand{\caveatfontfallback}{\itshape}
\IfFontExistsTF{Caveat}{%
  \newfontfamily\caveatfont{Caveat}[Scale=1.15]%
}{%
  \newcommand\caveatfont{\caveatfontfallback}%
}

\setmonofont{Latin Modern Mono}

% --- Mikro-Typografie ----------------------------------------------------
\usepackage{microtype}
\linespread{1.15}

% --- Farben & Boxen ------------------------------------------------------
\usepackage{xcolor}
\definecolor{lpInk}{RGB}{20,28,38}
\definecolor{kurs71}{HTML}{9D4A1A}
\definecolor{lpAccent}{HTML}{9D4A1A}
\definecolor{lpAccentLight}{RGB}{248,232,218}
\definecolor{lpAmber}{RGB}{222,138,30}
\definecolor{lpAmberLight}{RGB}{255,243,217}
\definecolor{lpAmberBorder}{RGB}{196,118,18}
\definecolor{lpGray}{RGB}{96,104,114}
\definecolor{lpGrayLight}{RGB}{238,240,243}

\usepackage[most]{tcolorbox}
\tcbuselibrary{breakable,skins}

\newtcolorbox{eselsbox}[1][Eselsbrücke]{%
  enhanced, breakable,
  colback=lpAmberLight,
  colframe=lpAmberBorder,
  boxrule=0.6pt,
  arc=2pt,
  borderline={0.4pt}{0pt}{lpAmberBorder, dashed},
  left=10pt,right=10pt,top=8pt,bottom=8pt,
  fonttitle=\sffamily\bfseries\color{lpAmberBorder},
  coltitle=lpAmberBorder,
  title={#1},
  attach boxed title to top left={xshift=8pt,yshift=-8pt},
  boxed title style={size=small, colback=lpAmberLight, colframe=lpAmberBorder, boxrule=0.4pt, arc=1pt},
  top=14pt
}

\newtcolorbox{merkbox}[1][Merksatz]{%
  enhanced, breakable,
  colback=lpAccentLight,
  colframe=lpAccent,
  boxrule=0.6pt,
  arc=2pt,
  left=10pt,right=10pt,top=8pt,bottom=8pt,
  fonttitle=\sffamily\bfseries\color{white},
  coltitle=white,
  title={#1},
  attach boxed title to top left={xshift=8pt,yshift=-8pt},
  boxed title style={size=small, colback=lpAccent, colframe=lpAccent, boxrule=0.4pt, arc=1pt},
  top=14pt
}

% --- Listen --------------------------------------------------------------
\usepackage{enumitem}
\setlist{itemsep=2pt, topsep=4pt, parsep=0pt}
\setlist[itemize,1]{label={\textbullet},leftmargin=1.6em}
\setlist[itemize,2]{label={\textendash},leftmargin=1.4em}
\setlist[enumerate,1]{label=\arabic*., leftmargin=1.8em}

% --- Tabellen ------------------------------------------------------------
\usepackage{tabularx}
\usepackage{array}
\usepackage{booktabs}
\usepackage{longtable}

% --- Kopf-/Fusszeilen ----------------------------------------------------
\usepackage{fancyhdr}
\usepackage{lastpage}
\pagestyle{fancy}
\fancyhf{}
\renewcommand{\headrulewidth}{0.3pt}
\renewcommand{\footrulewidth}{0pt}
\fancyhead[L]{\footnotesize\sffamily\color{lpGray}Lernplatt \textperiodcentered{} Kurs 71 \textperiodcentered{} Tag 7}
\fancyhead[R]{\footnotesize\sffamily\color{lpGray}Process Mining \textperiodcentered{} Discovery \textperiodcentered{} Conformance}
\fancyfoot[C]{\footnotesize\sffamily\color{lpGray}\thepage{} / \pageref{LastPage}}

\fancypagestyle{plain}{%
  \fancyhf{}%
  \renewcommand{\headrulewidth}{0pt}%
  \fancyfoot[C]{\footnotesize\sffamily\color{lpGray}\thepage{} / \pageref{LastPage}}%
}

% --- Section-Styling -----------------------------------------------------
\usepackage[explicit]{titlesec}

\titleformat{\section}[block]
  {\sffamily\Large\bfseries\color{lpAccent}}
  {\thesection.}{0.6em}{#1}
  [{\vspace{2pt}\color{lpAccent}\titlerule[0.6pt]}]

\titleformat{\subsection}[block]
  {\sffamily\large\bfseries\color{lpInk}}
  {\thesubsection}{0.6em}{#1}

\titleformat{\subsubsection}[block]
  {\sffamily\normalsize\bfseries\color{lpInk}}
  {}{0pt}{#1}

\titlespacing*{\section}{0pt}{14pt}{8pt}
\titlespacing*{\subsection}{0pt}{10pt}{4pt}
\titlespacing*{\subsubsection}{0pt}{8pt}{2pt}

% --- Hyperlinks ----------------------------------------------------------
\usepackage[hidelinks,unicode]{hyperref}

% --- TOC -----------------------------------------------------------------
\renewcommand{\contentsname}{\sffamily\Large\bfseries\color{lpAccent}Inhalt}

% --- Helfer --------------------------------------------------------------
\newcommand{\akzent}[1]{{\caveatfont\color{lpAmber}#1}}
\newcommand{\fachbegriff}[1]{\textbf{#1}}
\newcommand{\quelle}[1]{{\footnotesize\color{lpGray}(#1)}}

% =========================================================================
\begin{document}

% --- COVER ---------------------------------------------------------------
\begin{titlepage}
  \pagestyle{empty}
  \vspace*{1.5cm}
  \noindent{\sffamily\footnotesize\color{lpGray}LERNPLATT \textperiodcentered{} BY CAN SIEBERT}\\[2pt]
  \noindent{\color{lpAccent}\rule{\linewidth}{1.2pt}}\\[4pt]
  \noindent{\sffamily\footnotesize\color{lpGray}MODUL 3 \textperiodcentered{} TAG 7 VON 2 \textperiodcentered{} KURS 71 (Dig 42) \textperiodcentered{} 8.\,Juni 2026}

  \vspace{3cm}
  {\sffamily\fontsize{30}{34}\selectfont\bfseries\color{lpInk}Discovery\\\&\\Conformance\par}

  \vspace{0.7cm}
  {\sffamily\large\color{lpGray}Daten werden zur Realität --\\Vom Event-Log zum Ist-Prozess und zum Soll/Ist-Abgleich.\par}

  \vspace{1.5cm}
  {\caveatfont\LARGE\color{lpAmber}Lehrskript -- Tag 7 von 16}\par

  \vfill

  \noindent\begin{minipage}{0.55\linewidth}
    {\sffamily\footnotesize\color{lpGray}Lernpfad:}\\
    {\sffamily\small\color{lpInk}Wissen \textrightarrow{} Anwendung \textrightarrow{} Management}\\[10pt]
    {\sffamily\footnotesize\color{lpGray}Bearbeitungsdauer:}\\
    {\sffamily\small\color{lpInk}1 Kurstag (8 Unterrichtseinheiten)}
  \end{minipage}\hfill
  \begin{minipage}{0.4\linewidth}
    \raggedleft
    {\sffamily\footnotesize\color{lpGray}Dozent:}\\
    {\sffamily\small\color{lpInk}Can Siebert}\\[10pt]
    {\sffamily\footnotesize\color{lpGray}Format:}\\
    {\sffamily\small\color{lpInk}Theorie \textperiodcentered{} Fallstudie \textperiodcentered{} Coaching}
  \end{minipage}

  \vspace{0.5cm}
  \noindent{\color{lpAccent}\rule{\linewidth}{0.6pt}}
\end{titlepage}

% --- LERNZIELE -----------------------------------------------------------
\section*{Lernziele dieses Tages}
\addcontentsline{toc}{section}{Lernziele dieses Tages}
\thispagestyle{plain}

\noindent Modelle zeigen, wie ein Prozess \emph{gedacht} ist. Daten zeigen, wie er \emph{tatsächlich abläuft}. Diese Differenz ist die wichtigste Informationsquelle für Prozessverbesserung -- und gleichzeitig die häufigste Quelle für Audit-Findings. Heute führen wir die zwei zentralen Disziplinen des Process Mining ein: \fachbegriff{Discovery} (Ist-Prozess aus Event Logs ableiten) und \fachbegriff{Conformance Checking} (Abweichungen vom Soll erkennen, klassifizieren und behandeln).

\subsection*{L1 -- Wissen (Reproduktion und Einordnung)}
\begin{itemize}
  \item Event-Log-Mindestanforderungen kennen: Case ID, Aktivität, Timestamp (optional: Ressource, Kanal, Kosten, Standort).
  \item Discovery-Ergebnisarten benennen: Prozessgraph, Variantentabelle, Durchlaufzeiten, Engpasshypothesen.
  \item Conformance-Abweichungen klassifizieren: fehlende Schritte, zusätzliche Schritte, falsche Reihenfolge, Schleifen, SLA-Verletzungen.
  \item Den Unterschied zwischen ``Abweichung als Risiko'' und ``Abweichung als Innovation/Notfall'' verstehen.
\end{itemize}

\subsection*{L2 -- Anwendung (Transfer auf konkrete Situationen)}
\begin{itemize}
  \item Aus einem Mini-Event-Log Happy Path, Varianten und Engpässe ableiten -- ohne Tool.
  \item Abweichungen in Compliance-Risiko, Qualitäts-/Kostenrisiko, legitime Variante und Datenproblem einsortieren.
  \item Maßnahmen pro Abweichung definieren: Prozessregel, Training, Systemkontrolle oder Monitoring.
  \item Eine ``bewusst zugelassene Variante'' mit Guardrail formulieren.
\end{itemize}

\subsection*{L3 -- Management (Entscheidung und Verantwortung)}
\begin{itemize}
  \item Conformance als Steuerungsinstrument nutzen -- nicht als reines Reporting.
  \item Welche Abweichungen sind reputations-/compliancekritisch, welche Innovationspotenzial?
  \item Datenverantwortung benennen: Process Owner, Data Steward, Compliance Reviewer.
  \item Conformance-Programme unter Zielkonflikt (Speed vs.\ Compliance) verantworten.
\end{itemize}

\begin{merkbox}[Roter Faden]
Abweichungen vom Modell sind nicht automatisch Fehler. Sie sind \emph{Signale}: manche sind Compliance-Risiko, manche sind Innovation, manche sind Datenproblem. Die Senior-Disziplin liegt nicht im \emph{Finden} der Abweichungen -- sondern in der \emph{Klassifikation und der Konsequenz}.
\end{merkbox}

\clearpage

% --- 1. EINSTIEG ---------------------------------------------------------
\section{Einstieg: Vom Modell zur Realität}

In den vergangenen Tagen haben wir Modelle gezeichnet, in einem Repository steuerbar gemacht und über Workflows ausgeführt. Heute drehen wir den Spieß um: Wir beobachten, was \emph{tatsächlich} in den Systemen passiert, leiten daraus ein Ist-Bild ab, und vergleichen es mit unserem Soll-Modell. Diese Disziplin heißt \fachbegriff{Process Mining} und ist seit etwa 2010 von einem akademischen Spezialthema zu einer breit eingesetzten Management-Disziplin geworden. \quelle{van der Aalst, 2016; Reinkemeyer, 2020}

\subsection{Warum Daten oft eine andere Wahrheit erzählen}

Drei Beobachtungen aus Praxisprojekten:

\begin{enumerate}
  \item \fachbegriff{Modelle altern.} Ein BPMN-Diagramm von 2022 zeigt nicht den Prozess von 2026. Die Realität hat weitergedacht, das Modell ist eingefroren.
  \item \fachbegriff{Workarounds entstehen.} Wenn ein offizieller Schritt zu kompliziert ist, findet die Organisation einen Weg drumherum. Das ist nicht Sabotage, sondern Pragmatismus -- der das Modell zur Fiktion macht.
  \item \fachbegriff{Sonderfälle sind häufig.} Was im Modell als ``Ausnahme'' steht, ist in der Realität oft 30\,\% der Cases. Die Ausnahme ist die Regel.
\end{enumerate}

\subsection{Was Process Mining leistet}

\begin{itemize}
  \item \fachbegriff{Discovery.} Aus Event Logs einen Ist-Prozess rekonstruieren -- inklusive seltener Varianten, Schleifen, Durchlaufzeiten.
  \item \fachbegriff{Conformance Checking.} Den Ist-Prozess mit einem Soll-Modell vergleichen -- wo wird abgewichen, wie oft, mit welchem Effekt.
  \item \fachbegriff{Enhancement.} Modelle mit Daten anreichern -- Zeiten, Kosten, Ressourcen, Varianten (Thema von Tag 8).
\end{itemize}

\subsection{Die drei Bausteine eines Event Logs}

Ein \fachbegriff{Event Log} ist eine Tabelle mit Zeilen (Events) und Spalten (Eigenschaften). Drei Spalten sind Pflicht, ohne sie geht nichts: \quelle{IEEE, 2016; van der Aalst, 2012}

\begin{itemize}
  \item \fachbegriff{Case ID.} Identifiziert den konkreten Fall (Auftrag 4711, Ticket I-2026-0815). Alle Events mit derselben Case ID gehören zusammen.
  \item \fachbegriff{Aktivität (Activity).} Der Name des Schritts (``Auftrag angelegt'', ``Rechnung erstellt'', ``Versand''). Eindeutig benannt, idealerweise nach Verb + Objekt-Regel.
  \item \fachbegriff{Timestamp.} Wann ist das Event passiert? Mindestens Datum + Uhrzeit; bei kurzen Prozessen auch Millisekunden.
\end{itemize}

Optional, aber sehr nützlich:

\begin{itemize}
  \item \fachbegriff{Ressource.} Welcher Mitarbeiter, welches System hat den Schritt ausgeführt? (Datenschutz beachten.)
  \item \fachbegriff{Kanal.} Online, Filiale, Hotline, Partner?
  \item \fachbegriff{Kosten.} Hat dieser Schritt einen Wert in Aufwand oder \text{\euro}?
  \item \fachbegriff{Standort.} Welche Niederlassung, welches Land?
\end{itemize}

\begin{eselsbox}[Eselsbrücke: \akzent{C-A-T}]
Drei Pflichtspalten jedes Event Logs -- ohne sie kein Mining:\\[2pt]
\textbf{C}ase ID -- welcher Fall?\\
\textbf{A}ctivity -- welcher Schritt?\\
\textbf{T}imestamp -- wann?\\[2pt]
\emph{Wer C-A-T nicht in den Daten hat, kann Mining nicht starten -- nur Hoffen.}
\end{eselsbox}

\subsection{Typische Stolpersteine in Event Logs}

\begin{itemize}
  \item \fachbegriff{Fehlende oder unscharfe Activities.} Verschiedene Bezeichnungen für dieselbe Tätigkeit (``Approve'' vs.\ ``Genehmigen'').
  \item \fachbegriff{Doppelte Events.} Ein Klick erzeugt zwei Logzeilen, der Prozessgraph zeigt eine Schleife.
  \item \fachbegriff{Falsche Case IDs.} Ein Vorgang springt fälschlich zwischen zwei Cases.
  \item \fachbegriff{Zeitstempel-Probleme.} Zeitzonen, fehlende Sekunden, rückwärtslaufende Stempel durch Systemwechsel.
  \item \fachbegriff{Unvollständige Erfassung.} Manuelle Schritte werden nicht geloggt -- der Prozess ``springt''.
\end{itemize}

% --- 2. DISCOVERY --------------------------------------------------------
\section{Process Discovery -- vom Log zum Ist-Modell}

\fachbegriff{Process Discovery} ist die Disziplin, aus einem Event Log automatisch ein Prozessmodell zu erzeugen -- ohne dass jemand vorher ein Soll-Modell zeichnen muss. In der Praxis nutzen Tools (Celonis, UiPath Process Mining, Apromore, ProM, Disco u.\,a.) Algorithmen wie den \fachbegriff{Heuristics Miner}, den \fachbegriff{Inductive Miner} oder den \fachbegriff{Alpha-Algorithmus}, um aus Spuren typische Pfade und Varianten abzuleiten. \quelle{van der Aalst, 2016; van der Aalst \& Carmona, 2022}

\subsection{Was Discovery typischerweise liefert}

\begin{itemize}
  \item Einen \fachbegriff{Prozessgraph} mit Knoten (Aktivitäten) und Kanten (Übergänge), oft mit Häufigkeit oder Durchlaufzeit annotiert.
  \item Eine \fachbegriff{Variantentabelle} -- welche Pfad-Varianten kommen wie oft vor.
  \item \fachbegriff{Durchlaufzeiten je Schritt} und insgesamt, mit Quartilen.
  \item Erste \fachbegriff{Engpasshypothesen} -- wo stauen sich Cases?
\end{itemize}

\subsection{Discovery ohne Tool -- ein Mini-Log}

Wir machen den Schritt manuell, um die Logik zu verstehen. Vier Cases eines Auftragsprozesses:

\begin{quote}
\begin{tabular}{@{}lp{0.7\linewidth}@{}}
\toprule
Case 101 & Start \textrightarrow{} Anfrage\_prüfen \textrightarrow{} Angebot\_erstellen \textrightarrow{} Kunde\_bestätigt \textrightarrow{} Auftrag\_anlegen \textrightarrow{} Versand \textrightarrow{} Rechnung \textrightarrow{} Ende \\
Case 102 & Start \textrightarrow{} Anfrage\_prüfen \textrightarrow{} Rückfrage \textrightarrow{} Anfrage\_prüfen \textrightarrow{} Angebot\_erstellen \textrightarrow{} Kunde\_bestätigt \textrightarrow{} Auftrag\_anlegen \textrightarrow{} Rechnung \textrightarrow{} Versand \textrightarrow{} Ende \\
Case 103 & Start \textrightarrow{} Anfrage\_prüfen \textrightarrow{} Angebot\_erstellen \textrightarrow{} Kunde\_bestätigt \textrightarrow{} Auftrag\_anlegen \textrightarrow{} Versand \textrightarrow{} Versand \textrightarrow{} Rechnung \textrightarrow{} Ende \\
Case 104 & Start \textrightarrow{} Anfrage\_prüfen \textrightarrow{} Angebot\_erstellen \textrightarrow{} Kunde\_bestätigt \textrightarrow{} Auftrag\_anlegen \textrightarrow{} Storno \textrightarrow{} Ende \\
\bottomrule
\end{tabular}
\end{quote}

\subsubsection{Happy Path}

Der typischste Ablauf, der in den meisten Cases vorkommt:
\begin{quote}
Start \textrightarrow{} Anfrage\_prüfen \textrightarrow{} Angebot\_erstellen \textrightarrow{} Kunde\_bestätigt \textrightarrow{} Auftrag\_anlegen \textrightarrow{} Versand \textrightarrow{} Rechnung \textrightarrow{} Ende.
\end{quote}

(Case 101 entspricht dem Happy Path.)

\subsubsection{Drei Varianten}

\begin{itemize}
  \item \textbf{Variante A -- Rückfrage-Schleife (Case 102).} Ursache (Hypothese): Anfrage unvollständig, Daten fehlen. \emph{Risiko}: längere DLZ, mögliche Eskalation; \emph{möglicher Nutzen}: höhere Qualität der Angebote durch Klärung im Vorfeld.
  \item \textbf{Variante B -- Doppelversand (Case 103).} Ursache: möglicherweise getrennt versendete Teilmengen, oder Bedienfehler. \emph{Risiko}: zusätzliche Versandkosten, potenziell Kundenirritation; \emph{möglicher Nutzen}: bei großen Sendungen evtl.\ gewollt.
  \item \textbf{Variante C -- Storno nach Auftragsanlage (Case 104).} Ursache: Kundenrücktritt, Bonität, falsche Kalkulation. \emph{Risiko}: verlorene Bearbeitungszeit; \emph{möglicher Nutzen}: Lerneffekt, wenn häufig Stornos auftreten.
\end{itemize}

\subsubsection{Engpasshypothesen}

\begin{itemize}
  \item \fachbegriff{Anfrage\_prüfen} -- könnte Engpass sein, weil hier (Case 102) zurückgegangen wird und Rückfragen passieren.
  \item \fachbegriff{Übergang Auftrag\_anlegen $\rightarrow$ Versand} -- könnte Engpass sein, weil hier (Case 103) Doppelversand auftritt, was auf Koordinations\-probleme hindeutet.
\end{itemize}

\subsubsection{Drei Klärfragen an die Fachseite}

\begin{enumerate}
  \item Sind Rückfragen ein systematisches Problem (Datenqualität der Eingangsanfragen)?
  \item Ist der Doppelversand bewusst (Teillieferungen) oder ein Bedienfehler?
  \item Welcher Anteil der Cases endet im Storno, und welche Gründe sind dominant?
\end{enumerate}

\subsection{Was Discovery \emph{nicht} kann}

\begin{itemize}
  \item Discovery liefert \fachbegriff{Hypothesen}, keine Wahrheit. Jede gefundene Variante braucht Validierung mit der Fachseite.
  \item Discovery sieht nur, was \emph{geloggt} wird. Schritte, die manuell und ohne Log passieren, sind unsichtbar.
  \item Discovery ist datenqualitäts-abhängig. Schlechte Logs erzeugen Spaghetti-Grafen, die niemand versteht.
\end{itemize}

\begin{merkbox}[Discovery erzeugt Hypothesen, nicht Wahrheiten]
Ein Discovery-Modell sieht oft chaotisch aus. Das ist kein Tool-Problem -- es ist die Realität. Bevor man das Modell ``ordentlich'' macht, prüft die Senior-Disziplin: ist diese Variante echt, oder ein Datenartefakt? Ist diese Schleife Innovation, Fehler oder Pflicht? Erst nach Validierung wird aus dem Discovery-Graph ein belastbares Ist-Bild.
\end{merkbox}

\subsection{Filter-Techniken}

In realen Discovery-Projekten sind die Daten zu viel und zu unscharf, um sie direkt zu visualisieren. Drei nützliche Filter:

\begin{itemize}
  \item \fachbegriff{Frequenz-Filter.} Nur Pfade behalten, die in $\geq$X\,\% der Cases vorkommen.
  \item \fachbegriff{Zeitfenster-Filter.} Nur Cases aus einem bestimmten Zeitraum.
  \item \fachbegriff{Variant-Filter.} Nur die Top-N häufigsten Varianten anzeigen.
\end{itemize}

Diese Filter machen den Graphen lesbar -- aber sie verbergen Information, die später wieder hervorgeholt werden muss. Senior-Praxis: zweite Sicht ohne Filter, um Long-Tail-Phänomene nicht zu übersehen.

% --- 3. CONFORMANCE ------------------------------------------------------
\section{Conformance Checking -- Soll vs.\ Ist}

\fachbegriff{Conformance Checking} vergleicht den Ist-Prozess (aus Discovery) mit einem Soll-Modell (BPMN, EPC, Policy). Das Ergebnis sind \emph{Abweichungen} -- die zu klassifizieren und zu priorisieren sind. \quelle{Rozinat \& van der Aalst, 2008; van der Aalst, 2016}

\subsection{Vier typische Abweichungs-Arten}

\begin{enumerate}
  \item \fachbegriff{Fehlende Schritte.} Ein im Soll vorgesehener Schritt fehlt im Ist (z.\,B.\ Vier-Augen-Freigabe wurde umgangen).
  \item \fachbegriff{Zusätzliche Schritte.} Ein nicht vorgesehener Schritt erscheint (z.\,B.\ inoffizielle Sonderfreigabe).
  \item \fachbegriff{Falsche Reihenfolge.} Schritte in einer Reihenfolge, die das Modell nicht erlaubt (z.\,B.\ Rechnung vor Wareneingang).
  \item \fachbegriff{Schleifen oder Rework.} Schritte werden mehrfach durchgeführt (z.\,B.\ Anfrage\_prüfen wird wegen Rückfragen wiederholt).
\end{enumerate}

Hinzu kommen zwei spezielle Fälle:

\begin{itemize}
  \item \fachbegriff{SLA-Verletzungen.} Die Aktivitäten sind richtig, aber zu langsam.
  \item \fachbegriff{Datenprobleme.} Eine vermeintliche Abweichung beruht auf einem Log-Fehler, nicht auf einer realen Abweichung.
\end{itemize}

\subsection{Klassifikation -- die Senior-Disziplin}

Eine gefundene Abweichung ist erst dann steuerungstauglich, wenn sie in eine der vier Kategorien einsortiert wurde:

\begin{itemize}
  \item \fachbegriff{Compliance-Risiko.} Verstoß gegen Regel, Policy oder Gesetz. ``Zero Tolerance''-Kandidat.
  \item \fachbegriff{Qualitäts-/Kostenrisiko.} Operativ teuer oder qualitätsmindernd, aber nicht regulatorisch. Steuerbar über Prozessregeln, Training, Systemkontrolle.
  \item \fachbegriff{Legitime Variante.} Die Abweichung ist gewollt (Innovation, Service-Sonderfall) oder unvermeidlich. Ggf.\ ins Soll-Modell aufnehmen.
  \item \fachbegriff{Datenproblem.} Keine echte Abweichung, sondern Log-/Erfassungsfehler. Datenqualitäts-Maßnahme.
\end{itemize}

\subsection{Beispiel-Abweichungen aus dem Mini-Log}

\begin{tabularx}{\linewidth}{@{}lXl@{}}
\toprule
\textbf{Abweichung} & \textbf{Beschreibung} & \textbf{Kategorie} \\
\midrule
Rechnung vor Versand (Case 102) & Soll-Regel verletzt: ``Rechnung erst nach Versand'' & Compliance-Risiko \\
Doppelversand (Case 103) & ``Doppelversand unzulässig'' verletzt & Qualitätsrisiko (kann legitim sein, prüfen) \\
Rückfrage-Schleife (Case 102) & Modell erlaubt das, ok & Legitime Variante \\
Storno nach Auftragsanlage (Case 104) & Vor Versand erlaubt, ok & Legitime Variante \\
\bottomrule
\end{tabularx}

\medskip

\subsection{Priorisierung -- Schaden $\times$ Häufigkeit}

Wenn Dutzende Abweichungen identifiziert werden, braucht es eine Priorisierung. Eine pragmatische Heuristik: \emph{Schaden $\times$ Häufigkeit.} Schaden ist eine Schätzung (in \text{\euro}, in Compliance-Punkten, in Kundenwirkung); Häufigkeit ist messbar.

\begin{itemize}
  \item Hohe Häufigkeit, hoher Schaden $\rightarrow$ \emph{Sofort handeln}.
  \item Niedrige Häufigkeit, hoher Schaden $\rightarrow$ \emph{Kontrolle einbauen, Monitoring}.
  \item Hohe Häufigkeit, niedriger Schaden $\rightarrow$ \emph{Prozessoptimierung, Training}.
  \item Niedrige Häufigkeit, niedriger Schaden $\rightarrow$ \emph{tolerieren, beobachten}.
\end{itemize}

\subsection{Maßnahmen-Katalog}

Pro Abweichung gehört eine Maßnahme:

\begin{itemize}
  \item \fachbegriff{Prozessregel.} Anpassung der SOP, neue Pflichtfelder, klarerer Hauptpfad.
  \item \fachbegriff{Training.} Wissensvermittlung, Onboarding, Wiederholungen.
  \item \fachbegriff{Systemkontrolle.} Pflichtfeld, Validierungsregel, automatische Eskalation.
  \item \fachbegriff{Monitoring.} Dashboard, Alarmregel, regelmäßiger Review.
\end{itemize}

\subsection{Bewusst zugelassene Variante}

Nicht alles, was vom Soll abweicht, ist falsch. Manchmal ist eine Abweichung eine bewusste pragmatische Entscheidung -- die dokumentiert werden muss, sonst wird sie zum Audit-Befund. Beispiel:

\begin{itemize}
  \item Variante: ``Rechnung kann vor Wareneingang verbucht werden'' für Drop-Ship- oder Service-Leistungen.
  \item Guardrail: ``Nur für Lieferanten der Kategorie X, nur bis Betrag Y, mit Pflichtkennzeichen 'Drop-Ship'.''
  \item Nachweispflicht: monatlicher Audit-Report, dass alle Drop-Ship-Cases dieser Regel folgen.
\end{itemize}

\begin{eselsbox}[Eselsbrücke: \akzent{C-Q-L-D}]
Vier Kategorien für jede Conformance-Abweichung:\\[2pt]
\textbf{C}ompliance-Risiko -- Zero Tolerance.\\
\textbf{Q}ualitäts-/Kostenrisiko -- steuerbar, priorisierbar.\\
\textbf{L}egitime Variante -- ins Modell aufnehmen.\\
\textbf{D}atenproblem -- Log-/Erfassungsfehler beheben.\\[2pt]
\emph{Ohne Klassifikation ist jede Liste von Abweichungen ein Bermuda-Dreieck der Verantwortung.}
\end{eselsbox}

% --- 4. INTERPRETATION VS. AUTOMATIK -------------------------------------
\section{Interpretation: Abweichung als Signal, nicht als Fehler}

Eine zentrale Senior-Disziplin im Conformance-Geschäft ist die Haltung: \emph{eine Abweichung ist ein Signal, kein automatischer Befund.} \quelle{van der Aalst, 2012}

\subsection{Drei Arten, Signale zu lesen}

\begin{enumerate}
  \item \fachbegriff{Symptom-Lesart.} Abweichung zeigt, dass etwas nicht passt -- das Modell, der Prozess, die Daten, die Schulung. Frage: Was ist die Ursache?
  \item \fachbegriff{Innovations-Lesart.} Abweichung zeigt, dass die Realität ``vor'' dem Modell ist -- Mitarbeiter haben einen besseren Weg gefunden. Frage: Sollten wir den Weg ins Modell aufnehmen?
  \item \fachbegriff{Risiko-Lesart.} Abweichung zeigt, dass ein Risiko realisiert wird -- Compliance, Sicherheit, Qualität. Frage: Wie schließen wir die Lücke?
\end{enumerate}

\subsection{False-Compliance-Falle}

Eine besondere Gefahr ist die \fachbegriff{False Compliance}: Das Papier ist erfüllt, die Realität ignoriert. Beispiele:

\begin{itemize}
  \item Vier-Augen-Freigabe ``digital geklickt'', aber der zweite Klickende hat nicht reingeschaut.
  \item Kostenstellen-Pflichtfeld ist befüllt -- mit ``9999'' oder einer Catch-All-Stelle.
  \item Sonderfreigabe wird in 30\,\% der Cases gezogen, weil der Standardprozess zu langsam ist.
\end{itemize}

False Compliance ist kein Conformance-Fortschritt -- sie ist \emph{Verlagerung des Problems}. Senior-Disziplin: Conformance-Quoten nicht isoliert betrachten, sondern mit qualitativen Stichproben kombinieren.

\subsection{Zielkonflikt: Standardisierung vs.\ operative Flexibilität}

Wer Conformance zu streng durchsetzt, lähmt das operative Geschäft. Wer sie zu lasch durchsetzt, schafft Audit-Risiken. Praktische Lösung:

\begin{itemize}
  \item \textbf{Harte Regeln} für regulatorisch verpflichtete Aspekte (KYC, Datenschutz, Vier-Augen-Prinzip).
  \item \textbf{Weiche Regeln} mit klaren Eskalationspfaden für operative Flexibilität (Sonderfreigaben mit Schwellwert und Eskalations-Board).
  \item \textbf{Beobachtung} ohne automatische Konsequenz für Innovations-Potenziale (Pilot-Phase).
\end{itemize}

\begin{merkbox}[Conformance braucht Klassifikation]
Eine Abweichung ohne Klassifikation ist eine Suchanzeige ohne Verantwortlichen. Erst durch die Einsortierung in C-Q-L-D (Compliance, Qualität, Legitim, Daten) wird klar, wer handelt, mit welcher Dringlichkeit und mit welchem Werkzeug.
\end{merkbox}

% --- 5. FALLSTUDIE -------------------------------------------------------
\section{Fallstudie: MedTech Supplies -- Procure-to-Pay Conformance}

\emph{MedTech Supplies} ist ein reguliertes Unternehmen im Medizintechnik-Bereich. Auditoren bemängeln unklare Nachweise; gleichzeitig klagt der Einkauf über zu viel Bürokratie. Es gibt Hinweise auf ``Maverick Buying'' -- Einkäufe außerhalb des offiziellen Prozesses.

\subsection{Ausgangslage und Restriktionen}

\begin{itemize}
  \item Zeit: 6 Wochen bis Audit.
  \item Budget: 70\,000\,\text{\euro}.
  \item Logs: nur ERP und E-Mail-Tickets verfügbar -- unvollständig.
  \item Zielkonflikt: Speed (Einkauf) vs.\ Compliance (Audit).
\end{itemize}

\subsection{Schritt 1 -- Soll-Modell Procure-to-Pay (grob)}

\begin{enumerate}
  \item Bedarf erfassen (Anforderungsstelle).
  \item Bedarf genehmigen (Vorgesetzter, ggf.\ Budgetverantwortlicher).
  \item Bestellung anlegen (Einkauf).
  \item Lieferung erhalten.
  \item Wareneingang buchen (Lager).
  \item Rechnung erfassen (Finance).
  \item Drei-Wege-Abgleich (3-Way-Match: Bestellung, Wareneingang, Rechnung).
  \item Freigabe der Rechnung (Finance/Fachbereich).
  \item Zahlung auslösen (Finance).
  \item Archivierung / Audit Trail.
\end{enumerate}

\subsection{Schritt 2 -- Ist-Hinweise (qualitativ)}

\begin{itemize}
  \item 30\,\% Rechnungen kommen vor Wareneingang.
  \item 18\,\% haben fehlende Kostenstelle.
  \item 12\,\% laufen über ``Sonderfreigabe''.
  \item 8\,\% sind Dubletten.
\end{itemize}

\subsection{Schritt 3 -- Abweichungslandkarte}

\begin{tabularx}{\linewidth}{@{}>{\bfseries}lXll@{}}
\toprule
Abweichung & Risiko / Ursache (Hypothese) & Maßnahme & Owner \\
\midrule
Rechnung vor Wareneingang & 3-Way-Match bricht / Dropship, Service & Ausnahmeprozess, Pflichtkennzeichen, Nachbuchungs-SLA & Finance Ops \\
Fehlende Kostenstelle & Controlling, Compliance / schlechtes Formular & Pflichtfeld, Validierung, Rückgabe & Einkauf \\
Sonderfreigaben & Umgehung / Zeitdruck & Schwellwert, Eskalations-Board, Monitoring & Procurement Lead \\
Dubletten & Doppelzahlung / Datenqualität, Mehrfacherfassung & ID-Match, Sperrregel, Datenkonsolidierung & Finance Ops \\
\bottomrule
\end{tabularx}

\medskip

\subsection{Schritt 4 -- 6-Wochen-Plan (Audit-fähig)}

\begin{enumerate}
  \item Pflichtfelder Kostenstelle + Validierung im Anforderungsformular (Woche 1--2).
  \item Ausnahmeprozess für ``Rechnung vor Wareneingang'' formalisieren (Woche 2--3).
  \item Monitoring-Dashboard für Abweichungen + wöchentliches Review (Woche 3--6).
\end{enumerate}

\subsection{Schritt 5 -- 3-Monats-Plan (nachhaltig)}

\begin{enumerate}
  \item Automatisierter 3-Way-Match im ERP (Wochen 7--14).
  \item Lieferanten- und Stammdaten\-bereinigung (Wochen 8--16).
  \item Training für Anforderungsstellen + Governance-Update (Wochen 10--18).
\end{enumerate}

\subsection{Schritt 6 -- Bewusst zugelassene Variante}

\textbf{Variante:} ``Rechnung vor Wareneingang'' für Dropship-Service-Lieferanten.

\textbf{Guardrails:}
\begin{itemize}
  \item Nur für Lieferanten mit Markierung ``Dropship'' in Stammdaten.
  \item Nur bis Betrag 5\,000\,\text{\euro}.
  \item Pflichtkennzeichen ``Dropship'' im Rechnungs-Datensatz.
  \item Nachbuchungs-SLA: Wareneingang muss innerhalb 14\,Tagen erfolgen.
  \item Monatlicher Audit-Report mit Stichprobe.
\end{itemize}

\textbf{Frühwarnsignal:} Mehr als 5\,\% der Dropship-Cases ohne Nachbuchung innerhalb SLA $\rightarrow$ Regel re-evaluieren.

\subsection{Annahmen und Frühwarnsignale}

\begin{itemize}
  \item \textbf{Annahme:} Die Kostenstellen-Datenqualität verbessert sich durch Pflichtfeld + Validierung um $>$60\,\%.
  \item \textbf{Annahme:} Sonderfreigaben werden durch Schwellwert + Eskalation auf $<$5\,\% reduziert.
  \item \textbf{Frühwarnsignal 1:} Trotz Pflichtfeld bleiben Kostenstellen-Fehler $>$10\,\% $\rightarrow$ Validierung verschärfen.
  \item \textbf{Frühwarnsignal 2:} Sonderfreigaben steigen wieder $>$10\,\% $\rightarrow$ Eskalations-Board prüfen, Schwellwert anpassen.
\end{itemize}

\subsection{Lessons aus der Fallstudie}

\begin{itemize}
  \item Abweichungen müssen klassifiziert werden -- erst dann werden sie steuerbar.
  \item Eine bewusst zugelassene Variante ist Audit-fähig, wenn sie Guardrails und Frühwarnsignale hat.
  \item Speed-vs.-Compliance ist kein Entweder-Oder, sondern eine Frage der \emph{Schwellwerte und Eskalationspfade}.
\end{itemize}

% --- 6. DATENQUALITAET ---------------------------------------------------
\section{Datenqualität als Voraussetzung}

Discovery und Conformance liefern nur dann belastbare Ergebnisse, wenn die zugrunde liegenden Daten ``mining-tauglich'' sind. Die meisten Process-Mining-Projekte verbringen 40--60\,\% der Zeit mit Datenaufbereitung -- nicht mit Analyse. \quelle{van der Aalst, 2012; Reinkemeyer, 2020}

\subsection{Die fünf wichtigsten Datenqualitäts-Probleme}

\begin{enumerate}
  \item \fachbegriff{Inkonsistente Aktivitäten-Namen.} ``Approve'' / ``Genehmigen'' / ``Freigegeben'' werden als verschiedene Aktivitäten gezählt, obwohl es dieselbe ist. Maßnahme: Mapping-Tabelle, harmonisierte Naming.
  \item \fachbegriff{Lücken in Case IDs.} Zwei zusammengehörige Vorgänge tauchen unter verschiedenen IDs auf, oder ein Vorgang springt fälschlich zwischen IDs. Maßnahme: Case-ID-Strategie definieren (z.\,B.\ ``führende Order-ID'' für ganze P2P-Kette).
  \item \fachbegriff{Zeitstempel-Probleme.} Zeitzonen, fehlende Sekunden, rückwärtslaufende Stempel. Maßnahme: zentrale Normalisierung in Staging.
  \item \fachbegriff{Doppelte Events.} Ein Klick erzeugt zwei Logzeilen, der Graph zeigt Pseudo-Schleifen. Maßnahme: Deduplikation auf Basis Case+Activity+Timestamp (mit Toleranz).
  \item \fachbegriff{Manuelle Schritte ohne Log.} Aktivitäten, die per E-Mail oder Telefon passieren, fehlen vollständig. Maßnahme: Lücken explizit benennen, ggf.\ kurzfristige manuelle Erfassung.
\end{enumerate}

\subsection{Case-ID-Strategie}

Die zentrale Designentscheidung in jedem Mining-Projekt: \emph{Was ist die führende Case-ID?} Für P2P könnte das die Bestellnummer sein; für O2C die Auftragsnummer; für Incident-Management die Ticketnummer. Wichtig:

\begin{itemize}
  \item Pro Mining-Sicht gibt es \emph{eine} führende ID.
  \item Andere IDs (z.\,B.\ Rechnungsnummern) werden über Mapping mit der führenden ID verbunden.
  \item Wechsel der führenden ID erlaubt andere Sichten (z.\,B.\ Lieferanten-Sicht statt Bestell-Sicht).
\end{itemize}

\subsection{Datenqualitäts-SLAs}

Wer Process Mining im Dauerbetrieb fahren will, braucht \fachbegriff{Data Quality SLAs}: messbare Anforderungen an die Daten, mit Owner und Eskalationspfad.

\begin{itemize}
  \item \emph{Vollständigkeit Case ID:} $>$99\,\% der Events haben Case ID.
  \item \emph{Konsistenz Aktivitäten-Namen:} 100\,\% bekannt aus Mapping-Tabelle.
  \item \emph{Zeitstempel-Plausibilität:} 100\,\% innerhalb Quartal, keine Sprünge.
  \item \emph{Dubletten-Rate:} $<$0,5\,\% nach Deduplikation.
\end{itemize}

% --- 7. MANAGEMENT-PERSPEKTIVE -------------------------------------------
\section{Management-Perspektive: Conformance als Steuerungsinstrument}

Auf Wissens-Ebene sind Discovery und Conformance methodische Disziplinen. Auf Anwendungs-Ebene sind sie Handwerk. Auf Management-Ebene sind sie ein \emph{Steuerungs- und Compliance-Programm} -- mit Konsequenzen für Rollen, Datenschutz, Akzeptanz und Eskalationslogik.

\subsection{Zwei zentrale Zielkonflikte}

\begin{enumerate}
  \item \textbf{Compliance vs.\ Speed.} Strikte Conformance-Erzwingung erhöht Audit-Sicherheit -- und Bürokratie. Lockere Conformance bringt Tempo -- und Risiko. Empfehlung: \emph{harte Regeln für nicht verhandelbares, weiche Regeln mit Eskalation für operatives.}
  \item \textbf{Transparenz vs.\ Vertrauen.} Process Mining kann als ``Überwachung'' empfunden werden. Wer die Daten transparent macht, riskiert Akzeptanzverlust bei Mitarbeitenden. Empfehlung: Daten aggregiert nutzen, Personenbezug minimieren, Betriebsrat/Datenschutz früh einbinden.
\end{enumerate}

\subsection{Rollen im Conformance Operating Model}

\begin{itemize}
  \item \fachbegriff{Process Owner.} Fachliche End-to-End-Verantwortung. Entscheidet, welche Abweichungen toleriert werden.
  \item \fachbegriff{Data Steward.} Verantwortet Datenqualität, Mapping-Tabellen, Case-ID-Strategie.
  \item \fachbegriff{Compliance Reviewer.} Prüft regulatorische Aspekte, definiert Zero-Tolerance-Liste.
  \item \fachbegriff{Mining-Owner / Product Owner.} Verantwortet die Mining-Plattform als Produkt -- Releases, Use-Case-Pipeline, Stakeholder.
\end{itemize}

\subsection{Welche Abweichungen sind ``Zero Tolerance''?}

Eine kleine, aber harte Liste sollte jede Organisation explizit benennen. Typische Kandidaten:

\begin{itemize}
  \item Vier-Augen-Prinzip umgangen bei finanzwirksamen Freigaben.
  \item KYC-Schritte übersprungen oder unvollständig.
  \item Audit-Trail manipuliert oder gelöscht.
  \item Datenschutz-Rechte des Kunden umgangen (z.\,B.\ Löschanträge ignoriert).
\end{itemize}

\subsection{Datenlücken akzeptieren -- mit Kontrollen}

In der Realität sind nicht alle Datenquellen sofort minable. Senior-Praxis akzeptiert Lücken \emph{bewusst} und kompensiert sie:

\begin{itemize}
  \item Manuelle Schritte ohne Log: kompensieren durch wöchentliche manuelle Stichprobe.
  \item Lieferanten-IDs uneinheitlich: kompensieren durch Mapping-Tabelle und regelmäßiges Reconciliation.
  \item Zeitstempel-Lücken: kompensieren durch geschätzte Bandbreiten, mit Hinweis im Dashboard.
\end{itemize}

\subsection{Anti-Gaming der Conformance-KPIs}

Eine Fitness- oder Conformance-Quote als KPI hat denselben Gaming-Risiken wie jede andere. Beispiele:

\begin{tabularx}{\linewidth}{@{}lXX@{}}
\toprule
KPI & Gaming-Risiko & Gegenmaßnahme \\
\midrule
Conformance-Quote & Cases umetikettieren, damit sie ``passen'' & Stichproben mit qualitativem Audit \\
SLA-Hit-Rate & Cases früh ``vorläufig'' schließen & Verkoppeln mit Reopen-Rate \\
Abweichungs-Anzahl & Abweichungen nicht melden & Audit-Trail-Vergleich \\
\bottomrule
\end{tabularx}

\medskip

\subsection{Conformance als Steuerungsinstrument}

Wenn Process Mining nur Reporting erzeugt, bringt es wenig. Es entfaltet seinen Wert, wenn:

\begin{itemize}
  \item jede Abweichung eine \fachbegriff{Maßnahme} hat (Prozessregel, Training, Systemkontrolle, Monitoring).
  \item jede Maßnahme einen \fachbegriff{Owner} hat.
  \item jeder Owner eine \fachbegriff{Eskalationsregel} hat (was bei Unwirksamkeit der Maßnahme?).
  \item das Programm einen \fachbegriff{Lern-Rhythmus} hat (monatlicher oder vierteljährlicher Review).
\end{itemize}

\begin{merkbox}[Schluss-Merksatz für Tag 7]
Process Mining ohne Conformance ist Reporting; Conformance ohne Owner ist Liste; Owner ohne Eskalationsregel ist Theater. Erst die Kette \emph{Daten $\rightarrow$ Abweichung $\rightarrow$ Klassifikation $\rightarrow$ Maßnahme $\rightarrow$ Owner $\rightarrow$ Eskalation} macht aus Mining ein Steuerungssystem.
\end{merkbox}

% --- 8. TAKE-AWAYS -------------------------------------------------------
\section{Take-aways aus Tag 7}

\begin{enumerate}
  \item \textbf{Event Log = C-A-T.} Case ID, Activity, Timestamp -- ohne diese drei keine Discovery.
  \item \textbf{Discovery erzeugt Hypothesen, keine Wahrheiten.} Jede Variante mit Fachseite validieren.
  \item \textbf{Conformance braucht Klassifikation.} C-Q-L-D: Compliance, Qualität, Legitim, Daten.
  \item \textbf{Schaden $\times$ Häufigkeit} als pragmatische Priorisierungs-Heuristik.
  \item \textbf{Bewusst zugelassene Variante = Guardrails + Frühwarnsignal.} Sonst wird sie zum Audit-Befund.
  \item \textbf{False Compliance ist kein Fortschritt.} Stichproben mit qualitativer Tiefe kombinieren.
  \item \textbf{Datenqualität ist 50\,\% des Projekts.} Mapping, Case-ID-Strategie, Datenqualitäts-SLAs.
  \item \textbf{Mining-Programm braucht Operating Model.} Process Owner, Data Steward, Compliance Reviewer, Mining-Product-Owner.
\end{enumerate}

% --- 9. LITERATUR --------------------------------------------------------
\clearpage
\section{Literatur}

Im Skript verwendete und für die Vertiefung empfohlene Quellen. Zitierstil: APA-7.

\begin{description}[style=nextline,leftmargin=18pt,labelindent=0pt,itemsep=4pt]
  \item[Dumas, M., La Rosa, M., Mendling, J., \& Reijers, H.\,A. (2018).] \emph{Fundamentals of business process management} (2nd ed.). Springer. \href{https://doi.org/10.1007/978-3-662-56509-4}{https://doi.org/10.1007/978-3-662-56509-4}
  \item[IEEE Computational Intelligence Society. (2016).] \emph{IEEE Standard for eXtensible Event Stream (XES) for achieving interoperability in event logs and event streams} (IEEE Std 1849-2016). IEEE. \href{https://doi.org/10.1109/IEEESTD.2016.7740858}{https://doi.org/10.1109/IEEESTD.2016.7740858}
  \item[Reinkemeyer, L. (Hrsg.). (2020).] \emph{Process mining in action: Principles, use cases and outlook}. Springer. \href{https://doi.org/10.1007/978-3-030-40172-6}{https://doi.org/10.1007/978-3-030-40172-6}
  \item[Rozinat, A., \& van der Aalst, W.\,M.\,P. (2008).] Conformance checking of processes based on monitoring real behavior. \emph{Information Systems, 33}(1), 64--95. \href{https://doi.org/10.1016/j.is.2007.07.001}{https://doi.org/10.1016/j.is.2007.07.001}
  \item[van der Aalst, W.\,M.\,P. (2012).] Process mining manifesto. In F.~Daniel, K.~Barkaoui, \& S.~Dustdar (Hrsg.), \emph{Business process management workshops (BPM 2011)} (Lecture Notes in Business Information Processing, Vol.~99, S.~169--194). Springer. \href{https://doi.org/10.1007/978-3-642-28108-2_19}{https://doi.org/10.1007/978-3-642-28108-2\_19}
  \item[van der Aalst, W.\,M.\,P. (2016).] \emph{Process mining: Data science in action} (2nd ed.). Springer. \href{https://doi.org/10.1007/978-3-662-49851-4}{https://doi.org/10.1007/978-3-662-49851-4}
  \item[van der Aalst, W.\,M.\,P., \& Carmona, J. (Hrsg.). (2022).] \emph{Process mining handbook} (Lecture Notes in Business Information Processing, Vol.~448). Springer. \href{https://doi.org/10.1007/978-3-031-08848-3}{https://doi.org/10.1007/978-3-031-08848-3}
\end{description}

% --- 10. GLOSSAR ---------------------------------------------------------
\clearpage
\section{Glossar}

Die wichtigsten Begriffe des Tages -- kurz definiert, in alphabetischer Reihenfolge.

\begin{description}[style=nextline,leftmargin=0pt,labelindent=0pt,itemsep=4pt]
  \item[Activity (Aktivität)] Bezeichnung eines Schrittes im Event Log, üblich nach Verb + Objekt-Regel.
  \item[Alpha-Algorithmus] Klassischer Discovery-Algorithmus; Grundlage vieler Erweiterungen.
  \item[Case ID] Eindeutiger Identifier eines Vorgangs im Event Log. Bündelt alle zugehörigen Events.
  \item[Compliance-Risiko] Verstoß gegen Regel, Policy oder Gesetz; ``Zero Tolerance''-Kandidat.
  \item[Conformance Checking] Vergleich des Ist-Prozesses mit einem Soll-Modell zur Identifikation von Abweichungen.
  \item[Datenproblem (Abweichungs-Kategorie)] Vermeintliche Abweichung, die auf Log- oder Erfassungsfehler beruht.
  \item[Data Quality SLA] Messbare Anforderung an Eventlog-Daten, mit Owner und Eskalationspfad.
  \item[Data Steward] Rolle mit Verantwortung für Datenqualität, Mapping und Case-ID-Strategie.
  \item[Discovery (Process Discovery)] Disziplin, aus Event Logs einen Ist-Prozess automatisch abzuleiten.
  \item[Enhancement] Anreicherung eines Prozessmodells mit Daten (Zeiten, Kosten, Ressourcen). Thema Tag 8.
  \item[Event Log] Tabelle von Events mit mindestens Case ID, Activity und Timestamp.
  \item[False Compliance] Zustand, in dem die Form erfüllt ist, aber die Realität nicht (z.\,B.\ Vier-Augen-Klick ohne Prüfung).
  \item[Filter (Frequenz/Variant/Zeit)] Techniken, um Discovery-Graphen lesbar zu halten -- auf Kosten von Long-Tail-Information.
  \item[Fitness] Maß, wie gut ein Event Log zu einem Modell passt (in Conformance-Tools).
  \item[Guardrail] Bedingung, unter der eine bewusst zugelassene Variante akzeptabel ist.
  \item[Happy Path] Typischer Hauptpfad im Discovery-Graphen, in den meisten Cases vorkommend.
  \item[Heuristics Miner] Discovery-Algorithmus, der mit verrauschten Logs robust umgeht.
  \item[Inductive Miner] Discovery-Algorithmus mit garantiert ``sound'' erzeugten Prozessmodellen.
  \item[Legitime Variante] Abweichung, die nicht Fehler ist, sondern gewollt -- ggf.\ ins Modell aufnehmen.
  \item[Maverick Buying] Einkäufe außerhalb des offiziellen Prozesses; klassisches Conformance-Symptom.
  \item[Process Mining] Sammeldisziplin für Discovery, Conformance und Enhancement auf Basis von Event Logs.
  \item[Qualitätsrisiko] Abweichungs-Kategorie: operativ teuer oder qualitätsmindernd, aber nicht regulatorisch.
  \item[Schaden $\times$ Häufigkeit] Priorisierungs-Heuristik für Abweichungen.
  \item[Soll-Modell] Normatives Modell (BPMN, EPC, Policy), gegen das Conformance gemessen wird.
  \item[Variante (Process Mining)] Konkrete Pfad-Folge im Event Log, die sich vom Happy Path unterscheidet.
  \item[XES] \emph{eXtensible Event Stream.} IEEE-Standard für interoperable Event Logs.
\end{description}

\vspace{1cm}
\noindent{\color{lpAccent}\rule{\linewidth}{0.6pt}}\\[4pt]
{\footnotesize\sffamily\color{lpGray}Lernplatt \textperiodcentered{} by Can Siebert \textperiodcentered{} Kurs 71 (Dig 42) \textperiodcentered{} Tag 7 \textperiodcentered{} \textcopyright{} 2026}

\end{document}
